%%%%% Selecione a formatação para Publicação ou para Revisão %%%%%%%
%\documentclass[submissao]{workmicro2017} % Versão sem o nome dos autores
\documentclass[]{workmicro2017} % Versão com o nome dos autores

%\usepackage{showframe} %% mostra margens do documento. Util para edicao do texto. Comentar para copilacao da versao final

%%%%%%%%%%%%%%%%%%%


%%%%%%%% inclua suas opcoes pessoas (pacotes, comandos) aqui


%%%%%%%%%%%%%%%%%%%%%%%%%%%%%%%%%%%%%%%%

%%%%%%%%%%%%%%%%%%%%
%%% complete com os dados do trabalho %%%%%%%%
\titulo{Modelo de artigo para submissão no X Workshop Brasileiro de Micrometeorologia}
\autor[1]{João da Silva}
\autor[2]{Maria Antonieta Silveira}
\autor[3]{José Antunes}
\autor[2]{Joana Martins}
\autor[4]{Juarez da Silva}
\affil[1]{Curso de Graduação em Meteorologia, UFPEL}
\affil[2]{Curso de Graduação em Meteorologia, UFSC}
\affil[3]{Curso Técnico em Meteorologia, IFSC}
\affil[4]{Departamento de Física, UFSM}

\resumo{O resumo do trabalho deve ser escrito em um único parágrafo, com até 200 palavras. Ao final do resumo devem constar de 2 a 5 palavras-chave, separadas por ponto e vírgula.}
\palavrachave{palavra-chave 1; palavra-chave 2; palavra-chave 3}

\resumoingles{Para o resumo em língua inglesa (\textit{abstract}) se aplicam as mesmas regras do resumo em português.}
\keywords{keyword 1; keyword 2; keyword 3}

%%%%%%%%%%%%%%%%%%%%%%%%%%%%%%%%%%%%%%%%%%%%%%%%%%%%

\begin{document}





%\abstract{yyyy}

%%%%%%%%%% corpo do trabalho %%%%%%%%%%%%%%%%%%%%%%%%
\corpo  %%% nao remover!!!
%%%%%%%%%%%%%%%%%%%%%%%%%%%%%%%%%%%%%%
%%%% faca o seu trabalho a partir daqui!!!
%%%%%%%%%%%%%%%%%%%%%%%%%%%%


\section{Introdução}

\par O artigo submetido ao X Workshop Brasileiro de Meteorologia deve ser redigido de acordo com as seguintes normas: página a4 (norma ISO 216, com dimensões de $210mm\times297mm$), com orientação ``retrato'', contendo entre 3 e 6 páginas. As páginas devem ter margens de $2cm$ e numeradas a partir da segunda página. Cada parágrafo deve ser iniciado com um avanço (tabulação) de $125mm$ em relação à margem esqueda da página.

\par A numeração deve ser inserida de forma centralizada no rodapé da página, em fonte arial $10pt$, com um espaçamento de $0,5cm$ abaixo da margem inferior da página. Além disso, a partir da segunda página, no cabeçalho deve constar o seguinte texto, em negrito, fonte arial $9pt$, alinhado à direita: \textbf{X Workshop Brasileiro de Micrometeorologia, Santa Maria-RS, 8-10 de novembro de 2015.}. O texto do cabeçalho deve ter um espaçamento de $0,5cm$ em relação à margem superior da página.

\par O corpo principal do texto deve ser redigido em português ou inglês, com fonte arial em tamanho $11pt$ contendo, obrigatoriamente, resumos em lingua portuguesa e inglesa. Ainda, deve possuir as seções: introdução, materiais e métodos, resultados e conclusões (sendo os 3 últimos não necessariamente com estas nomenclaturas). Ao final o texto poderá usar um espaço para agradecimentos e, obrigatoriamente, uma seção de referências bibliográficas.



\section{Título do trabalho, autores, filiação e seções}

\par O título do trabalho deve ser apresentado centralizado no topo da primeira página, em fonte arial $14pt$ capitalizadas em tom de azul com a seguinte composição RGB: 31, 72, 118. Abaixo do título deve constar o nome do autor ou autora. Caso o trabalho possua mais de um autor ou autora os nomes devem ser separados por vírgula. Exceção se faz ao último nome que deve ser separado pela preposição ``e''. Os nomes devem ser inseridos com fonte arial $11pt$ e centralizados.

\par Cada autor ou autora deve ser associado à uma instituição. As filiações dos autores devem ser indicadas em fonte arial $10pt$, centralizadas e abaixo do nome dos autores, com uma linha de espaçamento. A ligação entre os nomes e filiações será realizada através de um algarismo arábico sobre-escrito logo após o nome do autor e imediatamente antes do nome da instituição.
\par Os títulos de seções devem ser inseridos sem indentação, em fonte arial $12pt$, negrito, na mesma cor do título. Para as seções numeradas, antes do nome da seção deve constar um algarísmo arábico, seguido por um ponto e com um espaço entre este ponto e o nome da seção. A numeração das seções deve seguir a ordem crescente a partir do número um. Entre a última linha de uma seção e o título da próxima deve haver uma linha de espaçamento.

\subsection{Subseções}
\par Caso seja necessário dividir uma seção em diferentes tópicos é permitido o uso de subseções. As subseções tem seu título formado de forma semelhante as seções porém são indentadas com o mesmo espaçamento de um parágrafo e sem o negrito. Quanto a numeração, as subseções devem ser iniciadas pelo número da seção na qual se inserem mais o número da subseção, separados por um ponto. A numeração da subseção também segue o mesmo critério das seções e deve ser reiniciado a cada seção.

\par O título da subseção deve possuir um espaçamento de uma linha em relação ao texto imediatamente anterior. 

\section{Apresentação das referências bibliográficas}

\par A referência a outros trabalhos acadêmicos deverão ser realizadas através do último sobrenome do autor e o ano de publicação. Tal referenciação deverá ser de forma indireta, como por exemplo: Segundo \citeonline{livro:stull1988}, a camada limite é caracterizada pelo escoamento turbulento..., ou ainda: a camada limite é caracterizada pelo escoamento turbulento \cite{livro:stull1988}. 

\par No caso de trabalhos com até três autores deve-se, no corpo do texto indicar os últimos sobrenomes seguido do ano de publicação: \cite{artigo:witek2011}. Trabalhos com mais de três autores devem ser apresentados pelo sobrenome do primeiro autor mais a expressão \textit{et al.}, seguido do ano de publicaçao: \cite{artigo:puhales2015}.

\par Recomenda-se os trabalhos referenciados sejam artigos publicados em periódicos, capítulos de livros ou livros na integra. No caso de capítulos, ou livros pode-se indicar a página específica da citação: \cite[p. 5]{livro:stull1988} ou \citeonline[p. 5]{livro:stull1988}. Caso seja necessário citar materiais de outras fontes exitem diversos exemplos na seção de referências bibliográficas.

\par Todos os documentos citados no texto devem ter sua referência expressa na seção de referências bibliográficas. Nas referências devem constar o nome de todos os autores, idependentemente da quantidade.


\section{Figuras, tabelas e equações}

\par As figuras, tabelas e equações apresentadas no trabalho devem ser numeradas, de forma crescente, a partir do número um. Cada um destes elementos tem sua numeração própria. De forma geral, tais elementos devem ser inseridos centralizados na página, com o espaçamento, superior e inferior, de uma linha em relação ao texto. Além disso, os elementos não devem ultrapassar as margens laterais do texto.

\subsection{Uso de figuras}

\par As figuras podem ser inseridas de três formas no texto. Uma delas, centralizada, dividindo o fluxo do texto, conforme o exemplo da figura \ref{fig1}.

\begin{figure}[ht!]
\centering
\includegraphics[width=0.65\textwidth]{figs/fig1.png}
\caption{\label{fig1}Representação da cascata de energia cinética turbulenta. Adaptada de \citeonline{artigo:puhales2015}.}
\end{figure}

\par Cada figura deverá ter uma legenda, posicionada abaixo da mesma com alinhamento justificado, em fonte arial $11pt$. No ínicio da legenda deve constar a palavra ``Figura'' seguida pelo seu respectivo número, como forma de identificação do elemento. Após esta identicação e separada por um travessão ``--'' deve constar uma descrição da figura, acompanhada da fonte se necessário.

\par Uma mesma figura pode ser dividida em subfiguras, e apresentada conforme a figura \ref{fig2}:

\begin{figure}[ht!]
\centering
\begin{subfigure}{0.45\textwidth}
\includegraphics[width=\textwidth]{figs/fig2a.png}
\caption{\label{fig2a}\centering Transição laminar-turbulento.}
\end{subfigure}
\begin{subfigure}{0.45\textwidth}
\includegraphics[width=\textwidth]{figs/fig2b.png}
\caption{\label{fig2b}\centering Taxa de dissipação normalizada.}
\end{subfigure}
\caption{\label{fig2}Características do escoamento turbulento. A figura \ref{fig2a} mostra a transição de um escoamento laminar para turbulento. Adaptada de \citeonline[p. 5]{cap:livro:wyngaard2010-c1}. A figura \ref{fig2b} apresenta o resultado encontrado neste estudo para a variação da taxa de dissipação com a velocidade do escoamento livre em um tunel de vento.}
\end{figure}

\par A disposição das sufiguras fica a critério do autor, respeitando-se as margens do texto. Além disso, cada subfigura deve ser identificado com uma letra, partindo-se da letra ``a'' e em ordem alfabética. Esta identificação deve ser feita logo abaixo da subfigura, entre parenteses e centralizada. Junto a cada subfigura deve constar uma descrição breve da mesma. Uma descrição mais completa deve ser inserida na legenda. Não há um limite para subfiguras de um mesmo elemento. Contudo, todos os painéis devem ficar na mesma página.

\begin{wrapfigure}[19]{r}{0.4\textwidth}
\begin{framed}
\centering
\includegraphics[width=\textwidth]{figs/fig3.png}
\caption{\label{fig3}Perfil vertical da taxa de dissipação.}
\end{framed}
\end{wrapfigure}
\par Adicionalmente, figuras podem ser inseridas lado a lado com o texto, conforme exemplifica a figura \ref{fig3}. Neste caso, a figura e sua legenda devem ser encerradas por um quadro. As considerações sobre a legenda e tamanho da figura são as mesmas.

\par Todas as figuras inseridas no texto devem ser referenciadas. Isto deve ser realizado através da palavra ``figura'' (com letras minúsculas), seguida do respectivo número (ou numero e letra no caso se subfigura). Ainda, não se faz exigências em relação à resolução mínima da figura. Contudo, se o comite avaliador julgar que a figura não está apresentada com qualidade gráfica aceitável, poderá exigir uma nova versão com maior qualidade.


\subsection{Uso de tabelas}

\par O uso de tabelas tem suas normas semelhantes àquelas para figuras que dividem o texto. Um exemplo é a tabela \ref{tab1}. Neste caso, faz-se diferença na legenda da tabela que deve ser apresentada no topo da mesma, iniciada pela palavra ``Tabela'' seguida de um algarísmo arábico e um travessão.

\begin{table}[h]
\centering
\caption{\label{tab1} Um exemplo de tabela. As tabelas também podem ser inseridas lado a lado com o texto, como a figura \ref{fig3}.}
\begin{tabular}{l l l l}
\hline
{col 1} & {col 2} & {col 3} & {col 4} \\
\hline
resl 1 & resl 2 & resl 3 & resl 4 \\
resl 11 & resl 12 & resl 13 & resl 14 \\
resl 21 & resl 22 & resl 23 & resl 24 \\
\hline
\end{tabular}
\end{table}

\par A tabela deve ter seu início e fim demarcados por uma linha horizontal. Além disso, após a primeira linha (de texto) da tabela deve constar uma outra linha horizontal. Recomenda-se tabelas curtas e que estejam contidas em apenas uma página. O alinhamento do texto nas colunas da tabela é livre. Todas as tabelas inseridas no texto devem ser referenciadas no mesmo através da palavra ``tabela'' mais o número que a identifica. Ainda, as tabelas também podem ser apresentadas lado a lado com o texto, semelhante ao apresentado para a figura \ref{fig3}.


\subsection{Uso de equações}

\par As equações devem ser apresentadas dividindo o fluxo do texto, como a equação \ref{eq1}. Além disso, deve ser apresentada centralizada e numerada por algarísmos arábicos em ordem crescente, a partir do número 1. A numeração das equações deve estar entre parenteses. Após a apresentação da equação, os símbolos ainda não referidos no texto devem ser claramente identificados:

\begin{equation}
\label{eq1}
\vec{F}=\frac{d\vec{p}}{dt}
\end{equation}

\noindent onde, na equação (\ref{eq1}), $\vec{F}$ é a força resultante, $\vec{p}$ é o \textit{momentum linear} e $t$  é o tempo. As equações também podem ser apresentadas na forma de subequações, como mostra os exemplos das equações (\ref{eq2a}) e (\ref{eq2b}):

\begin{subequations}
\begin{equation}
\label{eq2a}
\vec{F}=\frac{d}{dt}\left(m\vec{v}\right)
\end{equation}

\begin{equation}
\label{eq2b}
\vec{F}=m\frac{d\vec{v}}{dt}
\end{equation}
\end{subequations}

\par As equações numeradas devem ser referenciadas no texto através da palavra ``equação'' seguida pelo seu número identificador entre parenteses. No caso de subequações, segue-se o mesmo critério de numeração através de letras que o de subfiguras.

\par Caso seja necessário mostrar o desenvolvimento de uma formulação no trabalho, pode-se desenvolver a equação e numerar somente aquelas de interesse maior, que serão referenciadas no texto. Por exemplo:
{\allowdisplaybreaks
\begin{align}
\left<E\right>&=\frac{1}{2}\sigma_i^2\\
\left<E\right>&=\frac{1}{2}\left(\sigma_u^2+\sigma_v^2+\sigma_w^2\right)\nonumber\\
\left<E\right>&=\frac{1}{2}\sigma_u^2+\frac{1}{2}\sigma_v^2+\frac{1}{2}\sigma_w^2
\end{align}}
\noindent ou ainda:
\begin{subequations}
{\allowdisplaybreaks
\begin{align}
\vec{F}&=\frac{d\vec{p}}{dt}\nonumber\\
\vec{F}&=\frac{d}{dt}\left(m\vec{v}\right)\nonumber\\
\vec{F}&=m\frac{d\vec{v}}{dt}\\
\vec{F}&=m\vec{a}
\end{align}}
\end{subequations}

\par Assim como para figuras e tabelas, deve haver um espaçamento de uma linha entre as equações e o corpo do texto.

\section{Conclusão}

\par Os trabalhos deverão ser submetidos no formato PDF (\textit{portable document format}), seguindo as normas descritas neste modelo. Trabalhos submetidos fora das normas apresentadas serão  automaticamente rejeitados.
    
% % % % % % % % % % % % % % % % % % % % % % % % % % % % % % % % % % % % % % 	
% % % % % % % % % % % % % BIBLIOGRAFIA  % % % % % % % % % % % % % % % % % % 
% % % % % % % % % % % % % % % % % % % % % % % % % % % % % % % % % % % % % % 	

%%%%% BIBLIOGRAFIA -> INCLUIR NAS CHAVES O NOME DO ARQUIVO *.BIB
\bibliografia{referencias}
\nocite{evento:puhalesBLT2010,cap:livro:Nieuwstadt-VonDop:wyngaard81,livro:moninv1,livro:holton2004, tese:marquesfilho2004,diss:welter2006,tcc:mintegui2014,site:amsglo:fog}





\end{document}